\documentclass{report}
\usepackage{luatexja} % LuaTeXで日本語を使うためのパッケージ
\usepackage{luatexja-fontspec} % LuaTeX用の日本語フォント設定

% --- 数学関連 ---
\usepackage{amsmath, amssymb, amsfonts, mathtools, bm, amsthm} % 基本的な数学パッケージ
\usepackage{type1cm, upgreek} % 数式フォントとギリシャ文字
\usepackage{physics, mhchem} % 物理や化学の記号や式の表記を簡単にする

% --- 表関連 ---
\usepackage{multirow, longtable, tabularx, array, colortbl, dcolumn, diagbox} % 表のレイアウトを柔軟にする
\usepackage{tablefootnote, truthtable} % 表中に注釈を追加、真理値表
\usepackage{tabularray} % 高度な表組みレイアウト

% --- グラフィック関連 ---
\usepackage{tikz, graphicx} % 図の描画と画像の挿入
\usepackage{background} % ウォーターマークの設定
\usepackage{caption, subcaption} % 図や表のキャプション設定
\usepackage{float, here} % 図や表の位置指定

% --- レイアウトとページ設定 ---
\usepackage{fancyhdr} % ページヘッダー、フッター、余白の設定
\usepackage[top = 20truemm, bottom = 20truemm, left = 20truemm, right = 20truemm]{geometry}
\usepackage{fancybox, ascmac} % ボックスのデザイン

% --- 色とスタイル ---
\usepackage{xcolor, color, colortbl, tcolorbox} % 色とカラーボックス
\usepackage{listings, jvlisting} % コードの色付けとフォーマット

% --- 参考文献関連 ---
\usepackage{biblatex, usebib} % 参考文献の管理と挿入
\usepackage{url, hyperref} % URLとリンクの設定

% --- その他の便利なパッケージ ---
\usepackage{footmisc} % 脚注のカスタマイズ
\usepackage{multicol} % 複数段組
\usepackage{comment} % コメントアウトの拡張
\usepackage{siunitx} % 単位の表記
\usepackage{docmute}
% \usepackage{appendix}
% --- tcolorboxとtikzの設定 ---
\tcbuselibrary{theorems, breakable} % 定理のボックスと改ページ設定
\usetikzlibrary{decorations.markings, arrows.meta, calc} % tikzの装飾や矢印の設定

% --- 定理スタイルと数式設定 ---
\theoremstyle{definition} % 定義スタイル
\numberwithin{equation}{section} % 式番号をサブセクション単位でリセット

% --- hyperrefの設定 ---
\hypersetup{
  setpagesize = false,
  bookmarks = true,
  bookmarksdepth = tocdepth,
  bookmarksnumbered = true,
  colorlinks = false,
  pdftitle = {}, % PDFタイトル
  pdfsubject = {}, % PDFサブジェクト
  pdfauthor = {}, % PDF作者
  pdfkeywords = {} % PDFキーワード
}

% --- siunitxの設定 ---
\sisetup{
  table-format = 1.5, % 小数点以下の桁数
  table-number-alignment = center, % 数値の中央揃え
}

% --- 透かし画像の設定 --- 
\backgroundsetup{
  scale=0.5,                       % 画像のスケール
  % color=black,                   % 画像の色（透かし用に半透明が推奨）
  opacity=0.2,                   % 透かしの透明度（0が完全透明、1が完全不透明）
  angle=0,                       % 画像の角度
  position = current page.south east,  % ページの右下
  hshift=-6cm, % 右方向へのシフト（負の値で内側に移動）
  vshift=5cm,  % 上方向へのシフト（正の値で内側に移動）
  contents={\includegraphics{./fig/appilogo-circular-full.png}} % 画像のパス
}

% --- その他の設定 ---
\allowdisplaybreaks % 数式の途中改ページ許可
\newcolumntype{t}{!{\vrule width 0.1pt}} % 新しいカラムタイプ
\newcolumntype{b}{!{\vrule width 1.5pt}} % 太いカラム
\UseTblrLibrary{amsmath, booktabs, counter, diagbox, functional, hook, html, nameref, siunitx, varwidth, zref} % tabularrayのライブラリ
\setlength{\columnseprule}{0.4pt} % カラム区切り線の太さ
\captionsetup[figure]{font = bf} % 図のキャプションの太字設定
\captionsetup[table]{font = bf} % 表のキャプションの太字設定
\captionsetup[lstlisting]{font = bf} % コードのキャプションの太字設定
\captionsetup[subfigure]{font = bf, labelformat = simple} % サブ図のキャプション設定
\setcounter{secnumdepth}{4} % セクションの深さ設定
\newcolumntype{d}{D{.}{.}{5}} % 数値のカラム
\newcolumntype{M}[1]{>{\centering\arraybackslash}m{#1}} % センター揃えのカラム
\DeclareMathOperator{\diag}{diag}
\everymath{\displaystyle} % 数式のスタイル
\newcommand{\inner}[2]{\left\langle #1, #2 \right\rangle}
\renewcommand{\figurename}{図}
\renewcommand{\i}{\mathrm{i}} % 複素数単位i
\renewcommand{\laplacian}{\grad^2} % ラプラシアンの記号
\renewcommand{\thesubfigure}{(\alph{subfigure})} % サブ図の番号形式
\newcommand{\m}[3]{\multicolumn{#1}{#2}{#3}} % マルチカラムのショートカット
\renewcommand{\r}[1]{\mathrm{#1}} % mathrmのショートカット
\newcommand{\e}{\mathrm{e}} % 自然対数の底e
\newcommand{\Ef}{E_{\mathrm{F}}} % フェルミエネルギー
\renewcommand{\c}{\si{\degreeCelsius}} % 摂氏記号
\renewcommand{\d}{\r{d}} % d記号
\renewcommand{\t}[1]{\texttt{#1}} % タイプライタフォント
\newcommand{\kb}{k_{\mathrm{B}}} % ボルツマン定数
% \renewcommand{\phi}{\varphi} % ϕをφに変更
\renewcommand{\epsilon}{\varepsilon}
\newcommand{\fullref}[1]{\textbf{\ref{#1} \nameref{#1}}}
\newcommand{\reff}[1]{\textbf{図\ref{#1}}} % 図参照のショートカット
\newcommand{\reft}[1]{\textbf{表\ref{#1}}} % 表参照のショートカット
\newcommand{\refe}[1]{\textbf{式\eqref{#1}}} % 式参照のショートカット
\newcommand{\refp}[1]{\textbf{コード\ref{#1}}} % コード参照のショートカット
\renewcommand{\lstlistingname}{コード} % コードリストの名前
\renewcommand{\theequation}{\thesection.\arabic{equation}} % 式番号の形式
\renewcommand{\footrulewidth}{0.4pt} % フッターの線
\newcommand{\mar}[1]{\textcircled{\scriptsize #1}} % 丸囲み文字
\newcommand{\combination}[2]{{}_{#1} \mathrm{C}_{#2}} % 組み合わせ
\newcommand{\thline}{\noalign{\hrule height 0.1pt}} % 細い横線
\newcommand{\bhline}{\noalign{\hrule height 1.5pt}} % 太い横線

% --- カスタム色定義 ---
\definecolor{burgundy}{rgb}{0.5, 0.0, 0.13} % バーガンディ色
\definecolor{charcoal}{rgb}{0.21, 0.27, 0.31} % チャコール色
\definecolor{forest}{rgb}{0.0, 0.35, 0} % 森の緑色

% --- カスタム定理環境の定義 ---
\newtcbtheorem[number within = chapter]{myexc}{練習問題}{
  fonttitle = \gtfamily\sffamily\bfseries\upshape,
  colframe = forest,
  colback = forest!2!white,
  rightrule = 1pt,
  leftrule = 1pt,
  bottomrule = 2pt,
  colbacktitle = forest,
  theorem style = standard,
  breakable,
  arc = 0pt,
}{exc-ref}
\newtcbtheorem[number within = chapter]{myprop}{命題}{
  fonttitle = \gtfamily\sffamily\bfseries\upshape,
  colframe = blue!50!black,
  colback = blue!50!black!2!white,
  rightrule = 1pt,
  leftrule = 1pt,
  bottomrule = 2pt,
  colbacktitle = blue!50!black,
  theorem style = standard,
  breakable,
  arc = 0pt
}{proposition-ref}
\newtcbtheorem[number within = chapter]{myrem}{注意}{
  fonttitle = \gtfamily\sffamily\bfseries\upshape,
  colframe = yellow!20!black,
  colback = yellow!50,
  rightrule = 1pt,
  leftrule = 1pt,
  bottomrule = 2pt,
  colbacktitle = yellow!20!black,
  theorem style = standard,
  breakable,
  arc = 0pt
}{remark-ref}
\newtcbtheorem[number within = chapter]{myex}{例題}{
  fonttitle = \gtfamily\sffamily\bfseries\upshape,
  colframe = black,
  colback = white,
  rightrule = 1pt,
  leftrule = 1pt,
  bottomrule = 2pt,
  colbacktitle = black,
  theorem style = standard,
  breakable,
  arc = 0pt
}{example-ref}
\newtcbtheorem[number within = chapter]{exc}{Requirement}{myexc}{exc-ref}
\newcommand{\rqref}[1]{{\bfseries\sffamily 練習問題 \ref{exc-ref:#1}}}
\newtcbtheorem[number within = chapter]{definition}{Definition}{mydef}{definition-ref}
\newcommand{\dfref}[1]{{\bfseries\sffamily 定義 \ref{definition-ref:#1}}}
\newtcbtheorem[number within = chapter]{prop}{命題}{myprop}{proposition-ref}
\newcommand{\prref}[1]{{\bfseries\sffamily 命題 \ref{proposition-ref:#1}}}
\newtcbtheorem[number within = chapter]{rem}{注意}{myrem}{remark-ref}
\newcommand{\rmref}[1]{{\bfseries\sffamily 注意 \ref{remark-ref:#1}}}
\newtcbtheorem[number within = chapter]{ex}{例題}{myex}{example-ref}
\newcommand{\exref}[1]{{\bfseries\sffamily 例題 \ref{example-ref:#1}}}
% --- 再定義コマンド ---
% \mathtoolsset{showonlyrefs=true} % 必要な式番号のみ表示
\pagestyle{fancy} % ヘッダー・フッターのスタイル設定
\chead{応用量子物性講義ノート} % 中央ヘッダー
% \rhead{}
\fancyhead[R]{\rightmark}
\renewcommand{\sectionmark}[1]{\markright{\thesection\ #1}}
\cfoot{\thepage} % 中央フッターにページ番号
\lhead{}
\rfoot{Yuto Masuda and Haruki Aoki} % 右フッターに名前
\setcounter{tocdepth}{4} % 目次の深さ
\makeatletter
\@addtoreset{equation}{section} % サブセクションごとに式番号をリセット
\makeatother

% --- メタ情報 ---
\title{応用量子物性講義ノート}
\date{更新日\today}
\author{Yuto Masuda and Haruki Aoki}

\begin{document}
最後に，電磁場中の電子を考え，Dirac方程式の非相対論的極限($\abs{\bm{p}} \ll c$)がSchrödinger方程式であることを示す．
Dirac方程式は
\begin{align}
  \i\hbar\frac{\partial}{\partial t}\psi = (c \bm{\alpha}\cdot \bm{p} + \beta mc^2)\psi
\end{align}
である．ここに電磁場を次のように導入する．
\begin{align}
  \begin{dcases}
    \bm{p} \to \bm{p} + e\bm{A} (e>0)\\
    \phi = 0
  \end{dcases}
\end{align}
これを用いると
\begin{align}
  \i\hbar\frac{\partial}{\partial t}\psi = (c \bm{\alpha}\cdot (\bm{p} + e\bm{A}) + \beta mc^2)\psi \label{dirac-eq-in-emfield}
\end{align}
である．ここで，解の形を
\begin{align}
  \psi = \begin{pmatrix}
    \psi_1\\\psi_2\\\psi_3\\\psi_4
  \end{pmatrix}
  =
  \begin{pmatrix}
    \psi_a\\
    \psi_b
  \end{pmatrix}
\end{align}
とする．これを\refe{dirac-eq-in-emfield}に代入する．
\begin{align}
  \i\hbar\frac{\partial}{\partial t}
  \begin{pmatrix}
    \psi_a\\
    \psi_b
  \end{pmatrix}
  &=
  \begin{pmatrix}
    0 & c\bm{\sigma} \cdot (\bm{p} + e \bm{A})\\
    c\bm{\sigma} \cdot (\bm{p} + e \bm{A}) & 0
  \end{pmatrix}
  \begin{pmatrix}
    \psi_a\\
    \psi_b
  \end{pmatrix}
  +
  \begin{pmatrix}
    mc^2 & 0\\
    0 & -mc^2
  \end{pmatrix}
  \begin{pmatrix}
    \psi_a\\
    \psi_b
  \end{pmatrix}\\
  &=
  c\bm{\sigma}\cdot(\bm{p} + e\bm{A})
  \begin{pmatrix}
    \psi_b\\
    \psi_a
  \end{pmatrix}
  +mc^2
  \begin{pmatrix}
    \psi_a\\
    -\psi_b
  \end{pmatrix}
\end{align}
また，粒子のエネルギーを静止質量エネルギー($mc^2$)と非相対論的エネルギーの項($\epsilon_{\r{NR}}$)に分け，
\begin{align}
  \epsilon \sim mc^2 + \epsilon_{\r{NR}}
\end{align}
波動関数の時間発展を次のように記述する．
\begin{align}
  \begin{pmatrix}
    \psi_a\\
    \psi_b
  \end{pmatrix}
  =
  \begin{pmatrix}
    \psi_a^0\\
    \psi_b^0
  \end{pmatrix}
  \e^{-\i \frac{mc^2}{\hbar}t}
\end{align}
この表式を用いると\refe{dirac-eq-in-emfield}の左辺は
\begin{align}
  \i\hbar \frac{\partial}{\partial t}
  \begin{pmatrix}
    \psi_a\\
    \psi_b
  \end{pmatrix}
  =
  \i\hbar\qty[
    \frac{\partial}{\partial t}
    \begin{pmatrix}
      \psi_a^0\\
      \psi_b^0
    \end{pmatrix}
  ]
  \e^{-\i\frac{mc^2}{\hbar}t}
  + mc^2
  \begin{pmatrix}
    \psi_a^0\\
    \psi_b^0
  \end{pmatrix}
  \e^{-\i\frac{mc^2}{\hbar}t}
\end{align}
となる．これより，
\begin{align}
  \i\hbar\qty[
    \frac{\partial}{\partial t}
    \begin{pmatrix}
      \psi_a^0\\
      \psi_b^0
    \end{pmatrix}
  ]
  \e^{-\i\frac{mc^2}{\hbar}t}
  + mc^2
  \begin{pmatrix}
    \psi_a^0\\
    \psi_b^0
  \end{pmatrix}
  \e^{-\i\frac{mc^2}{\hbar}t}
  &=
  c\bm{\sigma}\cdot(\bm{p} + e\bm{A})
  \begin{pmatrix}
    \psi_b\\
    \psi_a
  \end{pmatrix}
  +mc^2
  \begin{pmatrix}
    \psi_a\\
    -\psi_b
  \end{pmatrix}\\
  \i\hbar\frac{\partial}{\partial t}
  \begin{pmatrix}
    \psi_a^0\\
    \psi_b^0
  \end{pmatrix}
  &=
  c\bm{\sigma}\cdot(\bm{p} + e\bm{A})
  \begin{pmatrix}
    \psi_b^0\\
    \psi_a^0
  \end{pmatrix}
  -2mc^2
  \begin{pmatrix}
    0\\
    \psi_b^0
  \end{pmatrix}
\end{align}
が得られる．上式の2行目に着目する．両辺を$mc^2$で割ると
\begin{align}
  \frac{1}{mc^2}\qty(\i\hbar\frac{\partial}{\partial t}\psi_b^0) = \frac{\bm{\sigma}\cdot(\bm{p} + e\bm{A})}{mc}\psi^0_a - 2 \psi^0_b
\end{align}
である．ここで左辺は
\begin{align}
  \frac{1}{mc^2}\qty(\i\hbar\frac{\partial}{\partial t}\e^{-\i\frac{\epsilon_{\r{NR}}}{\hbar} t}) \propto \frac{\epsilon_{\r{NR}}}{mc^2}
\end{align}
である．
また，
\begin{align}
  \epsilon_{\r{NR}} = \frac{\bm{p}^2}{2m}
\end{align}
だから非相対論的極限($\abs{\bm{p}} \ll c$)で左辺は0となる．よって
\begin{align}
  \psi_b^0 = \frac{\bm{\sigma}\cdot(\bm{p} + e\bm{A})}{2mc}\psi^0_a\\
\end{align}
を得る．これは粒子解と反粒子解を結ぶ式である．これを用いると1行目は
\begin{align}
  \i\hbar\frac{\partial}{\partial t}\psi^0_a &= c\bm{\sigma} \cdot (\bm{p} + e\bm{A})\psi_b^0\\
  &= \frac{[\bm{\sigma}\cdot(\bm{p} + e\bm{A})][\bm{\sigma}\cdot(\bm{p} + e\bm{A})]}{2m}\psi^0_a
\end{align}
と書き直される．次に$[\bm{\sigma}\cdot(\bm{p} + e\bm{A})][\bm{\sigma}\cdot(\bm{p} + e\bm{A})]$を計算する．
まず，Pauli行列には次の性質がある．
\begin{align}
  (\bm{\sigma}\cdot\hat{\bm{a}})(\bm{\sigma}\cdot\hat{\bm{b}}) &= (\sum_{i} \sigma_i \hat{a}_i) (\sum_{j} \sigma_j \hat{a}_j)\\
  &= \sum_{i} \sum_{j} \sigma_{i} \sigma_{j} \hat{a}_i \hat{b}_j\\
  &= \sum_{i} \sum_{j} \frac{1}{2}(\qty{\sigma_i, \sigma_j} + [\sigma_i, \sigma_j]) \hat{a}_i \hat{b}_j\\
  &= \sum_{i} \sum_{j} \frac{1}{2}(2 \delta_{i,j} + 2\i \sum_{k} \epsilon_{ijk} \sigma_k) \hat{a}_i \hat{b}_j\\
  &= \sum_{i} \hat{a}_i \hat{b}_i + \i \sigma_k \sum_{i} \sum_{j} \sum_{k} \epsilon_{ijk} \hat{a}_i \hat{b}_j\\
  &= \hat{\bm{a}}^2 + \i \sum_{k} \sigma_k (\hat{\bm{a}} \times \hat{\bm{b}})_k\\
  &= \hat{\bm{a}}^2 + \i \bm{\sigma} \cdot (\hat{\bm{a}} \times \hat{\bm{b}})
\end{align}
これを用いると
\begin{align}
  [\bm{\sigma}\cdot(\bm{p} + e\bm{A})][\bm{\sigma}\cdot(\bm{p} + e\bm{A})] &= (\hat{\bm{p}} + e\bm{A})^2 + \i \bm{\sigma} \cdot [(\hat{\bm{p}} + e\bm{A}) \times (\hat{\bm{p}} + e\bm{A})]\\
  &= (\hat{\bm{p}} + e\bm{A})^2 + \i e \bm{\sigma} \cdot [(\bm{A} \times \hat{\bm{p}}) + (\hat{\bm{p}} \times e\bm{A})]\\
\end{align}
と変形される．ここで，
\begin{align}
  [(\bm{A} \times \hat{\bm{p}}) + (\hat{\bm{p}} \times e\bm{A})]\psi &= (\bm{A} \times \hat{\bm{p}})\psi + (\hat{\bm{p}} \times \bm{A})\psi\\
  &= -\i \hbar \qty[(\bm{A} \times \nabla) \psi + (\nabla \times \bm{A}) \psi]\\
  &= -\i\hbar \qty[\bm{A} \times (\nabla \psi) + (\nabla\psi)\times \bm{A} + (\nabla \times \bm{A}) \psi]\\
  %&= \bm{A} \times (\hat{\bm{p}}\psi) + (\hat{\bm{p}}\psi) \times \bm{A} + (\hat{\bm{p}} \times \bm{A})\psi\\
  %&= (\hat{\bm{p}} \times \bm{A}) \psi\\
  &= -\i\hbar \qty[-(\nabla\psi)\times \bm{A} + (\nabla\psi)\times \bm{A} + (\nabla \times \bm{A}) \psi]\\
  &= -\i\hbar(\nabla \times \bm{A}) \psi
\end{align}
が成り立つ．以上より，
\begin{align}
  [\bm{\sigma}\cdot(\bm{p} + e\bm{A})][\bm{\sigma}\cdot(\bm{p} + e\bm{A})] &= (\hat{\bm{p}} + e\bm{A})^2 + e\hbar \bm{\sigma}\cdot(\nabla \times \bm{A})\\
  &= (\hat{\bm{p}} + e\bm{A})^2 + e\hbar \bm{\sigma} \cdot \bm{B}
\end{align}
したがって，
\begin{align}
  \i\hbar \frac{\partial}{\partial t} \psi_a^0 = \qty[\frac{1}{2m} (\hat{\bm{p}} + e\bm{A})^2 + \frac{e\hbar}{2m}\bm{\sigma}\cdot\bm{B}] \psi_a^0
\end{align}
が得られる．右辺第2項$\frac{e\hbar}{2m}\bm{\sigma}\cdot\bm{B}$はZeeman相互作用によるエネルギーを表す．これは磁場中の自由粒子のSchrödinger方程式と一致している．つまり，
電磁場中の自由粒子のDirac方程式の非相対論的極限は，電磁場中の自由粒子のSchrödinger方程式と一致する．
\end{document}